\chapter{Money}
\label{ch:money}

\section{What Ordenamiento got wrong}

Chapter~\ref{ch:crisis} described the currency unification of January 2021. This
chapter is written against it, so it is worth restating the failure precisely,
because the lesson is about sequence and not about direction.

Unification was right and had been right for twenty years. With the household
rate twenty-five-fold away from the enterprise rate, no Cuban firm's accounts
meant anything, and an economy whose accounts are fiction cannot be reformed
even by people who want to reform it. What went wrong was everything around it.

It was done with no reserves, and a fixed rate is a promise to sell dollars at
that rate. Cuba had no dollars, so the promise was not kept and a parallel market
reopened within weeks --- and once a parallel rate exists, the official rate
stops being a price and becomes an entitlement, allocated by whoever controls
access to it.

It was done with no supply response. Wages and pensions were multiplied several-
fold in an economy where the quantity of goods was falling. That is not a
distributional decision; it is the definition of an inflation.

It was done with the private sector still shut. The 2017 licence freeze was in
force and the law giving private firms legal personality was eight months away,
so the one part of the economy capable of responding to higher prices with more
output was legally prevented from doing so.

And it was done at the bottom of the worst external shock since 1993, when the
correct order is to stabilise first and unify afterwards.

The most damning comparison is internal. The same state, in 1994--96, with
fewer resources and a worse starting position, cut the deficit from around
thirty per cent of output to two, brought the informal exchange rate from a
hundred and fifty to twenty-five, and stabilised the peso without external
assistance --- and it did those things \emph{before} it did anything else.
Chapter~\ref{ch:dollarreform} tells that story. Cuba has demonstrated that it
knows how to do this in the right order. In 2021 it did not.

\section{The monetary regime}

The book takes a position here, and it is the most contestable recommendation in
Part III.

\emph{Cuba should officially and unilaterally dollarise, as its entry regime,
with a stated route out.}

The argument is not that dollarisation is optimal in general. It is not. It is
that the alternatives all require something Cuba does not have.

A \emph{managed float with an independent central bank} requires a central bank
whose commitments are believed. Cuba's has destroyed its citizens' savings twice
in thirty years --- in 1993--94 and again in 2021 --- and its ability to promise
anything about the value of money is, as a matter of observed fact, zero. Credibility
of that kind is rebuilt over a decade of demonstrated behaviour, and the design
needs a functioning currency in month one.

A \emph{currency board}, on the Baltic or Bulgarian model, keeps a national
currency fully backed by reserves. It is the right answer for a country that has
reserves. Cuba has none, cannot borrow them, and would have to acquire them
before the board could open --- which is a five-year detour at best.

A \emph{crawling peg} requires reserves to defend and discretion to adjust, and
discretion over an exchange rate in a country with Cuba's institutional history
is a rent, not a policy. Chapter~\ref{ch:integrity} will note that the single
largest corruption surface in the Cuban economy for thirty years has been access
to foreign currency at the official rate.

Dollarisation is the regime that requires no reserves, no credibility and no
discretion, which is exactly the combination Cuba's constraint set specifies.
It also happens to be where the country is already going: households price in
dollars, remittances arrive in dollars, and since 2024 the state's own shops
have been selling in dollars. The proposal is to make official and universal
what is already partial and inequitable --- because a partly dollarised economy
is the worst of all arrangements, dividing the population into those with access
to hard currency and those paid in something else, which is precisely the
inversion Chapter~\ref{ch:dollarreform} described and which is happening again.

The costs are real and must be stated. Cuba gives up monetary policy, gives up
the exchange rate as an adjustment mechanism, gives up seigniorage, and gives up
a lender of last resort. The first two matter less than they would elsewhere:
Cuba's export earners --- tourism, professional services, nickel, and the remote
services of Chapter~\ref{ch:talent} --- price in dollars already, so a
devaluation would not buy the competitiveness it buys a manufacturing exporter.
Seigniorage on a currency nobody wants to hold is small. The lender of last
resort is the genuine loss, and it binds hardest in exactly the scenario Cuba
faces regularly --- a Category 5 hurricane --- which is why
Design~\ref{d:rentrule}'s stabilisation fund and the contingent credit line
below are not optional extras but the other half of this regime choice.

Ecuador, El Salvador and Panama are the comparators. The Ecuadorean case is the
instructive one: dollarisation was adopted in 2000 in the middle of a banking
collapse, was expected by most economists to fail, and has survived twenty-five
years and several governments that disliked it --- because it turned out to be
enormously popular with people who had lost their savings, and no government has
been able to reverse it. That is a commitment device in the sense of
Chapter~\ref{ch:polity}: a policy that becomes hard to undo because the
population will not permit it.

\begin{design}{The monetary regime}
\label{d:regime}
\begin{rules}
\rl{1}{Adoption} The United States dollar is adopted as legal tender and unit of
account. The Cuban peso is retired on a published schedule at a single
conversion rate applied to all balances, wages, prices, contracts and public
accounts simultaneously.

\rl{2}{The conversion rate} Set at or close to the prevailing informal market
rate, not at an administrative rate. Setting it too strong is what destroys
savings; setting it at the market rate merely recognises a loss that has already
occurred, which is the difference between a reform and a confiscation.

\rl{3}{Single rate, immediately} There is no transitional dual rate, no
preferential rate for state enterprises, no allocation of currency at a
privileged rate to anybody, for any purpose, including food and fuel imports.
Every such arrangement in Cuban history has become a rent.

\rl{4}{Wage and pension protection} Public wages and pensions are converted at
the same rate and then re-set under Design~\ref{d:pay}, with a floor guaranteeing
that no pensioner's real income falls at conversion. This is the clause that
makes the regime survivable, and it is costed in Chapter~\ref{ch:sequence}.

\rl{5}{The exit} A published set of conditions --- reserves, fiscal record,
inflation, an established supervisory authority --- on which a national currency
may later be reintroduced, together with the procedure. Naming the exit is what
distinguishes a considered choice from a surrender, and it costs nothing to
write.
\end{rules}
\end{design}

\begin{design}{What replaces the central bank}
\label{d:bank}
\begin{rules}
\rl{1}{Function} Under dollarisation the central bank does not issue money. It
becomes a bank supervisor, a payments authority and the manager of the
stabilisation fund, with tenure, budget and reporting on the model of
Design~\ref{d:courts}.

\rl{2}{No fiscal financing} The prohibition on lending to the government is
absolute and constitutional. This is the single most important monetary clause in
the book, because monetary financing of the budget is how Cuba paid for its
welfare state in the years when neither Moscow nor Caracas was paying, and it is
what Chapter~\ref{ch:welfare} is designed to replace.

\rl{3}{Liquidity} A liquidity facility funded from the stabilisation fund and,
if obtainable, a contingent credit line with a multilateral or bilateral partner
--- the substitute for the lender of last resort that clause~1 gives up.
Obtaining it is a first-order objective of the debt normalisation below, not an
afterthought.

\rl{4}{Banking supervision} Capital, liquidity and disclosure standards on
international norms, applied to state and private banks identically, with a
resolution regime that exists before it is needed.
\end{rules}
\end{design}

\section{Debt, arrears and claims}

Cuba owes money to almost everyone and has not paid most of them. The stock is
several distinct problems that are conventionally handled separately and that
this book proposes to handle in one operation.

\emph{Paris Club official debt.} The 2015 arrangement forgave a large share of
accumulated arrears and rescheduled the remainder over eighteen years. Cuba
began missing payments within three years and the deal is in default.

\emph{Russian claims.} Ninety per cent of the Soviet-era debt was written off in
2014; the remainder was converted into investment commitments, largely
unrealised.

\emph{Chinese and other bilateral claims.} Substantially undisclosed, which is
itself a problem: a restructuring cannot be negotiated against an unknown
liability, and the first act of any process is to publish the stock.

\emph{Supplier arrears.} Unpaid invoices to foreign trading partners, in some
cases running to 2009. These are commercially the most damaging, because the
creditors are the same firms who would otherwise supply the imports the recovery
needs.

\emph{Frozen accounts.} Balances of foreign companies blocked from 2009 onward.
Legally these are not debt; commercially they are the reason nobody trusts a
Cuban account.

\emph{Expropriation claims.} The Foreign Claims Settlement Commission
registered 8,821 reports of seized assets and certified 5,913 of them, at a
principal of \$1.9 billion. At the 6 per cent simple interest the governing
statute specifies, that is now valued at something over \$9 billion
\citep{fcsc}. Only about
913 of the certified claims --- some 15 per cent --- are thought to meet the
conditions for a Title~III suit. Alongside them sits a far larger and
uncertified body of claims by Cuban Americans and by Cubans who never
left.

\begin{design}{One restructuring, not six}
\label{d:debt}
\begin{rules}
\rl{1}{Publish first} A complete, audited public statement of every external
obligation --- official, bilateral, commercial, arrears, frozen balances and
contingent --- before any negotiation opens. No restructuring has ever succeeded
against an undisclosed stock, and the disclosure is also the first act of the
transparency regime in Chapter~\ref{ch:integrity}.

\rl{2}{Single process} One comparable-treatment negotiation covering all classes
of claim, including the expropriation claims. Settling the sovereign debt while
leaving the expropriation claims outstanding leaves every foreign investor
exposed to Title III litigation and therefore leaves the country uninvestable;
settling the claims outside a general restructuring gives one class of creditor
priority and collapses the rest.

\rl{3}{Instruments} Long-dated, low-coupon bonds with a growth-linked component,
so that creditors participate in a recovery they are helping to finance and the
schedule does not bind in the years when the plan needs cash. Not a haircut
disguised as a maturity extension --- state the reduction and negotiate it.

\rl{4}{Claims settled in paper, within a cap} Per Chapter~\ref{ch:capital}, the
expropriation claims are settled by compensation, not restitution, in the same
instruments, within an aggregate cap that is announced in advance. No Cuban is
removed from a dwelling.

\rl{5}{Multilateral membership} Accession to the International Monetary Fund and
the multilateral development banks is pursued in parallel and is an objective of
the restructuring rather than a consequence of it, because their technical
assistance and their contingent lending are what clause~3 of
Design~\ref{d:bank} requires.

\rl{6}{Behaviour before terms} The frozen balances of Chapter~\ref{ch:venezuela}
are released and current supplier invoices are paid on time from the first
month, before and irrespective of any negotiation. A credit history is repaired
by conduct over years, not by a signed agreement, and this is the cheapest
credibility available.
\end{rules}
\end{design}

\section{The budget}

Chapter~\ref{ch:welfare} owns what is taxed. This chapter owns the deficit, how
it is financed, and the anchor that keeps the arrangement from unwinding.

\begin{design}{Fiscal architecture}
\label{d:fiscal}
\begin{rules}
\rl{1}{One budget} All state revenue and spending, including that of state
enterprises and of the military-linked conglomerates, appears in a single
published budget on international classification standards. At present a large
share of Cuban public finance is off-budget, which means nobody --- including
the government --- knows the fiscal position.

\rl{2}{No monetary financing} Per Design~\ref{d:bank}.2, and under dollarisation
it is mechanically impossible as well as prohibited, which is one of the
regime's principal attractions.

\rl{3}{Anchor} The recurrent balance rule of Design~\ref{d:rentrule}.3, with an
independent fiscal council reporting annually to the assembly and publishing its
own forecasts.

\rl{4}{Enterprise hard budgets} State enterprises face a hard budget constraint
with a published insolvency procedure. Subsidy to a loss-making enterprise, where
it continues, is an explicit and visible budget line, not a bank overdraft that
never gets called.

\rl{5}{Sequenced} The budget is brought into balance before the tax system is
capable of financing the welfare commitment, which means the first phase runs on
expenditure discipline and the rent rule, and the convergence to
Chapter~\ref{ch:welfare}'s benefit levels waits on the base. Saying which of
those comes first is the whole content of the word ``sequence''.
\end{rules}
\end{design}

\section{The ration book}

The \emph{libreta} was introduced in March 1962 as a temporary measure and is
sixty-four years old. It is a universal in-kind subsidy, badly targeted --- the
minister and the pensioner receive the same allocation --- expensive in foreign
exchange, and increasingly not delivered. It is also, for a substantial number of
elderly and eastern households, the difference between eating and not.

The transition is a genuine design problem rather than a matter of political
will, and getting it wrong starves people.

\begin{design}{From ration to income}
\label{d:libreta}
\begin{rules}
\rl{1}{Nothing is withdrawn before the replacement works} The universal cash
benefits of Design~\ref{d:benefit} must be paid, in full, reliably, through a
payment system that reaches every municipality, for a minimum of two
consecutive years before any ration item is removed.

\rl{2}{Replace, do not means-test} The libreta is converted into cash paid to
every household, not into a targeted benefit. Universality is retained for the
reason in Chapter~\ref{ch:welfare}: a means test is a discretionary decision and
Cuba cannot administer one honestly.

\rl{3}{Indexed} The cash amount is indexed to a published basket price, so that
the conversion is not a one-off transfer whose value inflation removes --- which
is what happened to the wage increases of January 2021.

\rl{4}{Sequence by item} Items are converted individually, beginning with those
whose market supply is adequate and ending with those whose supply depends on
Chapter~\ref{ch:land} succeeding. Rice, oil and milk go last.

\rl{5}{Reversible} If the replacement fails --- if supply collapses or the
payment system does not reach a province --- the ration is restored for that
item and that province, automatically, on a published trigger. This is the one
place in Part III where reversibility is a feature.
\end{rules}
\end{design}

\section{Payments}

The least glamorous section in Part III and among the most binding.

A tourist cannot pay in a Cuban restaurant with an ordinary card. A remote worker
in Havana cannot receive money from a client in Madrid into an account they
control. A private firm cannot pay a foreign supplier. A pensioner in Baracoa
cannot receive a transfer without a physical office. Chapter~\ref{ch:tourism}'s
value-per-visitor strategy, Chapter~\ref{ch:talent}'s entire proposition, and
Design~\ref{d:libreta}'s conversion all fail without payment rails, and none of
them can be built until the rails exist.

\begin{design}{Payment rails}
\label{d:pay2}
\begin{rules}
\rl{1}{Correspondent banking} Re-establishing correspondent relationships with
international banks is a first-order objective and the principal practical
reason for Design~\ref{d:debt}'s clause~6 and for the anti-money-laundering
standards in clause~4 below. Under the embargo this is difficult and, for
non-United States institutions, not impossible.

\rl{2}{Cards and acceptance} Card acceptance at any licensed business, with the
merchant fee capped and the settlement period fixed in regulation. A private
restaurant that can take a card is worth several times one that cannot.

\rl{3}{The right to hold and be paid} Any resident may hold a foreign-currency
account, receive international transfers into it, and spend from it without
administrative permission. This is the clause Chapter~\ref{ch:talent} depends on
entirely, and it is currently the binding constraint on Cuban service exports ---
more so than internet speed and far more so than visas.

\rl{4}{Standards} Anti-money-laundering and beneficial-ownership standards to
international norms, supervised by the authority in Design~\ref{d:bank}. These
are not a concession to foreigners; they are the price of the correspondent
relationships in clause~1, and they are the same register that
Chapter~\ref{ch:integrity} needs for domestic reasons.

\rl{5}{Reach} A basic transaction account available to every resident, including
by mobile, at every municipality. Financial inclusion is a precondition for
Design~\ref{d:libreta} and for the property tax of Design~\ref{d:tax}.4, and
without it both are administratively impossible in the east.
\end{rules}
\end{design}

\section{The first five hundred days}

The order, compressed. Chapter~\ref{ch:sequence} places this inside the whole
programme; here is the monetary sequence alone.

\emph{Days 1--100.} Publish the complete debt and arrears stock. Pay current
supplier invoices. Release the frozen balances. Announce the dollarisation date
and the conversion rate. Legislate the no-monetary-financing rule and the
recurrent balance rule. Open the debt negotiation and the IMF accession
conversation.

\emph{Days 100--250.} Convert. Single rate, all balances at once, wages and
pensions protected at the floor. Begin card acceptance and the basic account
rollout.

\emph{Days 250--500.} Stand up the revenue administration --- registration,
taxpayer numbers, withholding on the state payroll, the presumptive regime for
micro-enterprise --- knowing that it yields little for years. Begin the VAT
build. Establish the fiscal council and publish its first report. Begin the
libreta conversion on the two or three items whose market supply is adequate.

What is deliberately \emph{not} in the first five hundred days: raising benefit
levels to European standards, mass privatisation, and the withdrawal of any
ration item whose replacement is not already being paid. Each of those is a
later phase, and attempting them here is how the sequence fails.
